\documentclass[12pt]{article}
\usepackage[T1]{fontenc}
\usepackage[utf8]{inputenc}
\usepackage{lmodern}
\usepackage{textcomp}
\usepackage{enumitem}
\usepackage{geometry}
\usepackage{graphicx}
\usepackage{amsmath, amssymb}
\geometry{margin=1in}
\begin{document}
\begin{center}
Geography - Graduate Level Assessment 12 \
Total Marks: 40 \
Answer all questions.
\end{center}

\section*{Multiple Choice (10 marks)}
\begin{enumerate}[label=\arabic*.]
\item There are two billion children in the world today (2020), aged 0 to 15 years old. How many children will there be in the year 2100 according to the United Nations?
\begin{enumerate}[label=(\alph*)]
    \item 4 billion
    \item 3 billion
    \item 2 billion
    \item 1 billion
\end{enumerate}
\item As of 2017, what fraction of the population in Brazil used the internet in the past three months?
\begin{enumerate}[label=(\alph*)]
    \item 18\%
    \item 38\%
    \item 58\%
    \item 78\%
\end{enumerate}
\item As of 2015, agriculture made up about what percentage of total Indian GDP?
\begin{enumerate}[label=(\alph*)]
    \item 8\%
    \item 16\%
    \item 32\%
    \item 64\%
\end{enumerate}
\item As of 2019, about what percentage of Russians say it is very important to have free media in our country without government/state censorship?
\begin{enumerate}[label=(\alph*)]
    \item 38\%
    \item 53\%
    \item 68\%
    \item 83\%
\end{enumerate}
\item What was GDP per capita in the United States in 1950 when adjusting for inflation and PPP in 2011 prices?
\begin{enumerate}[label=(\alph*)]
    \item About \$3 k
    \item About \$8 k
    \item About \$15 k
    \item About \$30 k
\end{enumerate}
\end{enumerate}

\section*{True / False (10 marks)}
\textit{Answer True or False}
\begin{enumerate}[label=\arabic*.]
\setcounter{enumi}{5}
\item True or False: Frederick lost much of his artillery but retreated in good order, helped by dense woods. The Austrians had ultimately made little progress in the campaign in Saxony despite Hochkirch and had failed to achieve a decisive breakthrough. After a thwarted attempt to take Dresdena different answer, so that Saxony remained under Prussian occupation.
\item True or False: Carthage was Palermo's major trading partner under the Phoenicians and the city enjoyed a prolonged peace during this period. Palermo came into contact with the Ancient Greeks between the 6 th and the 5 th centuries BC which preceded the a different answer, a conflict fought between the Greeks of Syracuse and the Phoenicians of Carthage for control over the island of Sicily.
\item True or False: But early in the 16 th century, a brilliant new red appeared in Europe. When the Spanish conquistador a different answer and his soldiers conquered the Aztec Empire in 1519-1521, they discovered slowly that the Aztecs had another treasure beside silver and gold; they had the tiny cochineal, a parasitic scale insect which lived on cactus plants, which, when dried and crushed, made a magnificent red.
\item True or False: Theatre traditionally has occupied a large place in Parisian culture, and many of its most popular actors today are also stars of French television. The oldest and most famous Paris theatre is the Comédie-Française, founded in 1680.
\item True or False: The practice spread to include generals, admirals, politicians, doctors and scientists such as Isaac Newton, buried on 4 April 1727, and Charles Darwin, buried 26 April 1882. Another was a different answer who led the movement to abolish slavery in the United Kingdom and the Plantations, buried on 3 August 1833.
\end{enumerate}

\section*{Long Form Response (20 marks)}
\textit{Answer each question with a clear and complete explanation.}
\begin{enumerate}[label=\arabic*.]
\setcounter{enumi}{10}
\item Read the following passage and answer the question: On 2 March 2009, however, Vieira was assassinated by what preliminary reports indicated to be a group of soldiers avenging the death of the head of joint chiefs of staff, General Batista Tagme Na Wai. Tagme died in an explosion on Sunday, 1 March 2009, target of an assassination. Military leaders in the country pledged to respect the constitutional order of succession. National Assembly Speaker Raimundo Pereira was appointed as an interim president until a nationwide election on 28 June 2009. It was won by Malam Bacai Sanhá. Question: Who won the election in June 2009?
\item Read the following passage and answer the question: On October 17, 1987, about 3,000 Armenians demonstrated in Yerevan complaining about the condition of Lake Sevan, the Nairit chemicals plant, and the Metsamor Nuclear Power Plant, and air pollution in Yerevan. Police tried to prevent the protest but took no action to stop it once the march was underway. The demonstration was led by Armenian writers such as Silva Kaputikian, Zori Balayan, and Maro Margarian and leaders from the National Survival organization. The march originated at the Opera Plaza after speakers, mainly intellectuals, addressed the crowd. Question: How many people demonstrated?
\end{enumerate}

\vfill
\noindent\textit{End of Paper}
\end{document}